\chapter{Проверка результата}

\section{Формулировка и аргумент}
\begin{enumerate}
\item Все объекты определены; порядок кванторов зафиксирован.
\item Гипотезы проверены на примере и устанавливаются без предполагаемого
      вывода.
\item У конечного входа указано представление чисел и структур;
      у оценки --- параметры и модель вычислений.
\item Для каждой леммы выписаны зависимости и проверены условия внешних теорем.
\item Различены свободная группа и группа точного действия,
      стабилизатор уровня и стабилизатор одной вершины.
\item Согласованы сторона действия, порядок произведения секций,
      переход к фактору и независимость от базиса.
\item Ранг вычислен точной арифметикой; конечные значения
      и утверждение для всех глубин имеют отдельные доказательства.
\end{enumerate}
Перечитать аргумент от заключения к предпосылкам. Затем изложить ключевую
лемму без рабочих записей и проверить удаление одной гипотезы.
К вычислениям приложить точные входные данные.

\section{Новизна и внешняя проверка}
Для ближайшей теоремы записать номер, страницу, гипотезы, вывод и отношение
к кандидату. Проверить нужный фрагмент доказательства, предшественников
и продолжения. На неделе~2 подготовить вопрос для обсуждения;
на неделях~4 и~8 --- конкретную лемму с доказательством.
Для замечания записать ответ и изменение текста.

Первая теорема закончена при наличии полного доказательства,
проверенного примера, сравнения с литературой и связной записки.
Внешний отзыв фиксируется отдельно с датой и объёмом проверки.
До его получения результат имеет статус авторской проверки.
Решение о публикации принимается после оценки вклада и разбора замечаний.
